\documentclass[10pt,a4paper]{amsart}
\usepackage[margin=19mm]{geometry}
\usepackage{amsmath,amssymb,mathtools}
\usepackage[scr=rsfs,cal=euler]{mathalfa}
\usepackage{xfrac}
\usepackage[unicode,hidelinks,bookmarks=false]{hyperref}
\newcommand{\ie}{i.e.\ }
\newcommand{\cf}{cf.\ }
\newcommand{\resp}{respectively\ }
\newcommand{\Subsection}{\subsection}
\providecommand{\nobreakdash}{\nobreak\hbox{-}}
\setlength{\emergencystretch}{2em}
\multlinegap0pt
\usepackage{bm}
\usepackage{bbm}
\usepackage{dsfont}
\usepackage{xspace}
\usepackage{cancel}
\usepackage{mathtools}
\usepackage{amsthm}
\makeatletter
\@ifundefined{theorem}{\newtheorem{theorem}{Theorem}}{}
\@ifundefined{lemma}{\newtheorem{lemma}[theorem]{Lemma}}{}
\makeatother
\def\N{{\mathbb N}} \def\bbO{{\mathbb O}}
\newcommand{\R}{\mathbb{R}}
\newcommand{\tilv}{\tilde{v}}
\title[Gamma-Convergence of Higher-Order Phase Transition Models]{Gamma-Convergence of Higher-Order Phase Transition Models}
\author{Denis Brazke, Gianna G\"otzmann, Hans Kn\"upfer}
\date{Source: arXiv:2503.08309v1 (CC BY 4.0)}
\begin{document}
\maketitle

\noindent\emph{Source and licence.} Gamma-Convergence of Higher-Order Phase Transition Models. Version arXiv:2503.08309v1 (CC BY 4.0), distributed under the Creative Commons Attribution 4.0 International licence, as recorded in the arXiv licence metadata for this exact version and shown on the abstract page https://arxiv.org/abs/2503.08309v1. Full rights notice: licenses/arxiv-2503.08309-CC-BY-4.0.txt.\par\smallskip

\noindent\textit{Editorial scope.} Counted unit: the Lemma whose source label is \texttt{IntLemAbstr} in the article (source0.tex lines 1205--1210, source label \texttt{IntLemAbstr}), delivered together with the article's own complete original proof of it (source0.tex lines 1212--1250, one \texttt{proof} environment of 39 lines). No mathematics is written, repaired or re-derived: statement and proof are the article's own text line for line. Required context retained uncounted: none. Every definition, display and label of those lines is retained unchanged. Omitted from this package and not redistributed: 1--1204, 1211--1211, 1251--1290. Nothing inside a retained excerpt is omitted or altered: every retained body is delivered exactly as the article's lines are written, so no omission receipt of any kind accompanies this package. Citations: the retained excerpts carry no citation command at all, so no bibliography entry of the article is transcribed and no reference is redistributed. Reference adaptation: none was needed and none was made; every reference inside the delivered unit resolves inside it, so no unresolved reference or citation is produced. Printed slips kept literal: no printed word, symbol, hypothesis or number of the counted statement or of its proof was altered. Layout and class: the article's own class is \texttt{article} with packages including inputenc, xcolor, latexsym, amssymb, amsmath, amsthm, amsxtra, amsbsy, accents, mathtools, mathalpha, bm, hyperref, subfig; the wrapper is the lane's pinned \texttt{amsart} editorial wrapper at 10pt on A4, which restates the theorem-like environments the retained text needs and adds the article's own macro definitions that the retained text needs (\texttt{\textbackslash N}, \texttt{\textbackslash R}, \texttt{\textbackslash tilv}), with extra wrapper packages (bm, bbm, dsfont, xspace, cancel, mathtools). The delivered numbering is the unit's own. Assets: no retained span contains a figure environment, a graphics-inclusion command, a PDF-page inclusion, a table or a TikZ picture; no image file is copied into this package, the package declares no asset dependency and no image file is redistributed with it. Licence: version arXiv:2503.08309v1 of this article is distributed under the Creative Commons Attribution 4.0 International licence, as recorded in the arXiv licence metadata for this exact version and shown on the abstract page https://arxiv.org/abs/2503.08309v1. A full rights notice is delivered at licenses/arxiv-2503.08309-CC-BY-4.0.txt. Characters: every retained excerpt is pure ASCII, so the delivered engine, XeLaTeX with its default Latin Modern text font, reports no missing character.
\begin{lemma} \label{IntLemAbstr}
Let $I \subset \R$ be an open bounded interval, $1 \leq q,r < \infty$ and $j,m \in \N_0$ where $j<m$. Then, there exists $C=C(q,r,m,I)>0$ such that for all $u \in L^q(I) \cap W^{m,r}(I)$, we have
\begin{align} \label{eq:IntLemAbstr}
    \|u^{(j)}\|_{L^r(I)} \leq \|u^{(m)}\|_{L^r(I)} + C\|u\|_{L^q(I)}.
\end{align}
\end{lemma}

\begin{proof}
We argue by contradiction and suppose there exists a sequence $u_n \in W^{m,r}(I)$ fulfilling
    \begin{align} \label{assume}
        \|u_n^{(m)}\|_{L^r(I)} + n\|u_n\|_{L^q(I)} \leq \|u_n^{(j)}\|_{L^r(I)}.
    \end{align}
We define the sequence $v_n \coloneq u_n \left(\|u_n\|_{W^{m-1,r}(I)}+\|u_n^{(j)}\|_{L^r(I)}\right)^{-1} \in W^{m,r}(I)$ such that
\begin{align} \label{constseq}
    \|v_n\|_{W^{m-1,r}(I)}+\|v_n^{(j)}\|_{L^r(I)}=1.
\end{align}
Since \eqref{assume} also holds for the sequence $v_n$, we then obtain
\begin{align} \label{bounded1}
    \|v_n^{(m)}\|_{L^r(I)}+n\|v_n\|_{L^q(I)} 
    \leq \|v_n^{(j)}\|_{L^r(I)}
    \overset{\eqref{constseq}}{\leq} 1-\|v_n\|_{W^{m-1,r}(I)} \leq 1,
\end{align}
and thereby
\begin{align} \label{bounded2}
\begin{split}
    \|v_n\|_{W^{m,r}(I)} & \overset{\eqref{bounded1}}{\leq}  \|v_n\|_{W^{m-1,r}(I)} +1 - n \|v_n\|_{L^q(I)}\\
    &\overset{\eqref{constseq}}{\leq} 2- \|v_n^{(j)}\|_{L^r(I)} - n \|v_n\|_{L^q(I)}
    \leq 2.
\end{split}
\end{align}
Hence, there exist $\tilv \in W^{m,r}(I)$ and a subsequence of $v_n$ (not relabeled), such that $v_n^{(i)} \rightharpoonup \tilv^{(i)}$ in $L^r(I)$ as $n \to \infty$ for all $i \in \{0,\ldots , m\}$. Since $W^{1,r}(I)$ is compactly embedded into $L^r(I)$, we obtain strong convergence of a further subsequence for the derivatives up to order $m-1$, that is
\begin{align}
    \|v_n-\tilv \|_{W^{m-1,r}(I)}+ \|v_n^{(j)}- \tilv^{(j)}\|_{L^r(I)} = \sum_{i=0}^{m-1} \|v_n^{(i)}- \tilv^{(i)}\|_{L^r(I)} +\|v_n^{(j)}- \tilv^{(j)}\|_{L^r(I)} \to 0
\end{align}
as $n \to \infty$. At the same time, we have 
\begin{align}
    \|v_n-\tilv \|_{W^{m-1,r}(I)}+ \|v_n^{(j)}- \tilv^{(j)}\|_{L^r(I)} 
    \overset{\eqref{constseq}}{\geq} 1 - \|\tilv \|_{W^{m-1,r}(I)} - \|\tilv^{(j)}\|_{L^r(I)},
\end{align}
implying $\|\tilv \|_{W^{m-1,r}(I)} + \|\tilv^{(j)}\|_{L^r(I)} \geq 1$.
From \eqref{bounded1}, we further deduce
\begin{align}
    \lim_{n\to \infty} \|v_n\|_{L^q(I)} \leq \lim_{n \to \infty} \frac 1n \left(1-\|v_n^{(m)}\|_{L^r(I)}\right) \overset{\eqref{bounded2}}{=} 0.
\end{align}
Thus, $v_n$ has a subsequence that converges to zero almost everywhere. On the other hand, we find a further subsequence that converges almost everywhere to $\tilv$ due to the strong convergence in $L^r(I)$. Consequently, we have $\tilv = 0$ almost everywhere, which is a contradiction.
\end{proof}

\end{document}
